%==============================================================================
\section{Phân cụm dựa trên centroid (Centroid-Based Clustering)}
\label{sec:centroid-based}
%==============================================================================

\begin{dinhnghia}[title={Phân cụm dựa trên centroid}]
Phân cụm dựa trên centroid là lớp thuật toán học máy nhằm \textbf{phân hoạch} tập dữ liệu thành các nhóm (cụm) dựa trên \textbf{khoảng cách} từ điểm dữ liệu đến centroid của từng cụm.
\end{dinhnghia}

\begin{ynghia}[title={Centroid là gì?}]
Centroid của một cụm là \textbf{trung bình cộng số học} của mọi điểm trong cụm đó; nó đóng vai trò điểm đại diện cho cụm.
\end{ynghia}

\begin{phuongphap}[title={Hai thuật toán phổ biến nhất}]
Hai thuật toán phân cụm dựa trên centroid được dùng rộng rãi nhất là \textbf{K-means} và \textbf{K-medoids}.
\end{phuongphap}

\subsection{K-means Clustering}

\begin{dinhnghia}[title={K-means}]
K-means là thuật toán \textbf{học không giám sát} phổ biến để phân cụm dữ liệu: đơn giản, hiệu quả, gom điểm vào $K$ cụm theo độ tương tự.
\end{dinhnghia}

\begin{phuongphap}[title={Cơ chế tóm tắt}]
\begin{enumerate}
  \item Chọn ngẫu nhiên $K$ centroid ban đầu --- tâm khởi tạo của mỗi cụm.
  \item Gán mỗi điểm vào cụm có centroid \textbf{gần nhất}.
  \item Cập nhật centroid bằng \textbf{trung bình} các điểm trong cụm.
  \item Lặp đến khi centroid không đổi hoặc đạt số vòng lặp tối đa.
\end{enumerate}
\end{phuongphap}

\subsection{K-Medoids Clustering}

\begin{dinhnghia}[title={K-medoids}]
K-medoids là thuật toán phân hoạch gom tập điểm thành $k$ cụm. Khác K-means dùng \textbf{giá trị trung bình} làm tâm cụm, K-medoids dùng một \textbf{điểm dữ liệu thực} gọi là \textbf{medoid} làm đại diện.
\end{dinhnghia}

\begin{ynghia}[title={Medoid}]
Medoid là điểm trong cụm \textbf{tối thiểu hóa tổng khoảng cách} tới mọi điểm khác trong cụm đó.
\end{ynghia}

\begin{tinhchat}[title={So với K-means}]
K-medoids \textbf{bền hơn} với ngoại lai và nhiễu vì tâm cụm luôn là một mẫu quan sát được, không phải trung bình có thể bị kéo lệch bởi điểm cực đoan.
\end{tinhchat}
